\hypertarget{appendix-3c73f7862157}{}
\section*{Exceptional projections and dimension}
\section{Exceptional projections and dimension}\label{exceptional-projections-and-dimension}

\textbf{Curriculum depth:} Developed · \textbf{First seen:} University courses · \textbf{Latest development:} University courses

\subsection{Definition and mechanism}\label{definition-and-mechanism}

Projection usually lowers dimension, but a finite or fractal set may collapse unusually many points in special directions. For a finite planar set \(X\), count pairs whose orthogonal projections coincide in direction \(\theta\); such pairs lie on lines parallel to the projection fibers. Averaging over directions gives a baseline, while exceptional directions produce larger fibers and encode geometric structure. For a density, the related transform integrates along each fiber, as in tomography; the Fourier-slice identity connects the transform of a projection to a slice of the original Fourier transform. Dimension and separation assumptions matter: finite point sets, thickened sets, measures, and finite-field analogues have different statements. A highly structured set can defeat naive random-grid intuition, and subfield examples in finite fields create strong obstructions. Results are estimates with hypotheses, not universal formulas for every dimension or projection family. This node connects geometric projection to incidence counting and harmonic analysis.

\subsection{Prerequisites}\label{prerequisites}

\begin{itemize}
\tightlist
\item
  \href{}{Basis coordinates and dimension (01-Foundations-I, appendix-ce0341f7ea0d)}
\item
  \href{}{Projection rejection and scalar components (02-Foundations-II, appendix-505a502c372a)}
\item
  \hyperref[appendix-b60686ba06a1]{Fourier analysis and spectral phase}
\item
  \href{}{Vector spaces subspaces and span (02-Foundations-II, appendix-173b52c063ae)}
\end{itemize}

\subsection{Used by}\label{used-by}

\begin{itemize}
\tightlist
\item
  \href{Continuous\%20Notes/University\%20-\%20Worlds,\%20Images\%20and\%20Inference\%20-\%20Graphics,\%20Vision\%20and\%20Photography\#Finite-field\%20projections\%20turn\%20image\%20size\%20into\%20collision\%20counts}{Finite-field projections turn image size into collision counts}
\item
  \href{Continuous\%20Notes/University\%20-\%20Worlds,\%20Images\%20and\%20Inference\%20-\%20Graphics,\%20Vision\%20and\%20Photography\#Real\%20projections\%20add\%20scale,\%20spacing,\%20and\%20incidence\%20geometry}{Real projections add scale, spacing, and incidence geometry}
\item
  \hyperref[appendix-b0732b22e7d6]{Additive energy and approximate structure}
\end{itemize}

\subsection{Full teaching explanation}\label{full-teaching-explanation}

\begin{itemize}
\tightlist
\item
  \href{Continuous\%20Notes/University\%20-\%20Worlds,\%20Images\%20and\%20Inference\%20-\%20Graphics,\%20Vision\%20and\%20Photography\#Finite-field\%20projections\%20turn\%20image\%20size\%20into\%20collision\%20counts}{Finite-field projections turn image size into collision counts}
\item
  \href{Continuous\%20Notes/University\%20-\%20Worlds,\%20Images\%20and\%20Inference\%20-\%20Graphics,\%20Vision\%20and\%20Photography\#Real\%20projections\%20add\%20scale,\%20spacing,\%20and\%20incidence\%20geometry}{Real projections add scale, spacing, and incidence geometry}
\end{itemize}

\subsection{Recall and next depth}\label{recall-and-next-depth}

\begin{itemize}
\tightlist
\item
  Relate a coincident projected pair to a line fiber and a direction.
\item
  Why do a finite set, a thickened set, and a finite-field set require different hypotheses?
\end{itemize}

\subsection{Source and navigation}\label{source-and-navigation}

Teaching narrative: \href{Continuous\%20Notes/University\%20-\%20Worlds,\%20Images\%20and\%20Inference\%20-\%20Graphics,\%20Vision\%20and\%20Photography\#Finite-field\%20projections\%20turn\%20image\%20size\%20into\%20collision\%20counts}{open the main explanation}. Browse the \href{}{University graphics and vision index (02-Foundations-II, appendix-a6d556a4ef50 § University graphics vision and photography - independent series)}. Additional chapter connections: \href{Continuous\%20Notes/University\%20-\%20Worlds,\%20Images\%20and\%20Inference\%20-\%20Graphics,\%20Vision\%20and\%20Photography\#Cameras,\%20projection,\%20and\%20calibration}{Cameras, projection, and calibration}; \href{Continuous\%20Notes/University\%20-\%20Worlds,\%20Images\%20and\%20Inference\%20-\%20Graphics,\%20Vision\%20and\%20Photography\#Sampling\%20continuous\%20signals\%20and\%20image\%20formation}{Sampling continuous signals and image formation}; \href{Continuous\%20Notes/Computing\%20Foundations\%20-\%20From\%20Bits\%20to\%20C++\%20and\%20Systems\%20APIs\#Projection\%20separates\%20useful\%20movement\%20from\%20sideways\%20movement}{Projection separates useful movement from sideways movement}.

Caption/source records: \href{Continuous\%20Notes/University\%20-\%20Worlds,\%20Images\%20and\%20Inference\%20-\%20Graphics,\%20Vision\%20and\%20Photography\#P2-01\%20-\%20Lecture\%2001:\%20Introduction\%20to\%20Projection\%20Theory}{P2-01}, \href{Continuous\%20Notes/University\%20-\%20Worlds,\%20Images\%20and\%20Inference\%20-\%20Graphics,\%20Vision\%20and\%20Photography\#P2-02\%20-\%20Lecture\%2002:\%20Fundamental\%20Methods\%20of\%20Projection\%20Theory}{P2-02}, \href{Continuous\%20Notes/University\%20-\%20Worlds,\%20Images\%20and\%20Inference\%20-\%20Graphics,\%20Vision\%20and\%20Photography\#P2-09\%20-\%20Lecture\%2009:\%20Reflections\%20on\%20the\%20Szemeredi-Trotter\%20Theorem}{P2-09}, \href{Continuous\%20Notes/University\%20-\%20Worlds,\%20Images\%20and\%20Inference\%20-\%20Graphics,\%20Vision\%20and\%20Photography\#P2-17\%20-\%20Lecture\%2017:\%20Random\%20Walks\%20on\%20Finite\%20Groups,\%20Part\%201}{P2-17}.
